

\subsection{TODO into 1: When Bayes Became Computable: Monte Carlo, Metropolis--Hastings, and Gibbs Sampling}
\label{sec:mcmc-history}

During World War II, Alan Turing and I. J. Good demonstrated that Bayesian updating could be operationalised through sequential evidence weighting—using log-odds, \emph{decibans}, and Good--Turing smoothing to solve cryptanalytic problems under strict mechanical constraints. 

However, these wartime methods relied either on discrete evidence accumulation or continuous analytical approximations. In the post-war decades, as Bayesian inference expanded into complex, high-dimensional parameter spaces $\Theta \subseteq \mathbb{R}^d$, calculating posterior expectations $\mathbb{E}_\pi[h(\theta)] = \int_\Theta h(\theta)\pi(\theta \mid y)d\theta$ required evaluating the marginal likelihood (normalising constant):

\textcolor{red}{Reference to Bayesian into subsection once written up}

\begin{equation}
    p(y) = \int_{\Theta} p(y \mid \theta)\pi_{0}(\theta)\,d\theta.
    \label{eq:marginal-likelihood}
\end{equation}
In high dimensions, standard numerical quadrature scales as $\mathcal{O}(K^d)$ for a grid of size $K$, rendering exact evaluation computationally intractable - the modern problem of the \emph{curse of dimensionality}.

While we will detail a small subsection of the history and methods developed, more details about the history of the MCMC methods can be found in \cite{robert2011short, grimmett2014probability, wiki:mcmc, wiki:metropolishastings}


\noindent
\fcolorbox{red}{red!10}{%
  \parbox{\dimexpr\linewidth-2\fboxsep-2\fboxrule\relax}{%
    \raggedright
    \textcolor{red}{%
      \textbf{TODO:} Add
      
      \begin{enumerate}
        \item overall examples of where mcmc methods are used today - locating the flight?
        \item maybe more details about what inference is vs bayesian inference and why we need it? -> what do 18yo students actually know about this?
        \item explaining in simple terms the consequences of the theorems + their connection to CLT --> elaborate further + add proofs from books + mention Markov Chains intro 1040s
        \item Links to the full algorithms in appendix or code on github. 
        \item any images of the mcmc e.g. locality theorem of gibbs. 
        \item link to Markov Chain CLT with metropolis hastings algo --> using it
        \item What is a transition kernel density and what are the conditions for it? 
      \end{enumerate}
    }%
  }%
}

% -----------------------------------------------------------------------------
\subsubsection{Physics Origins: The Metropolis--Hastings Framework (1953--1970)}
\label{subsec:metropolis-hastings}

The computational breakthrough behind modern Markov chain Monte Carlo (MCMC) originated not in statistics, but in physical chemistry. Nicholas Metropolis, Arianna Rosenbluth, Marshall Rosenbluth, Augusta Teller, and Edward Teller \cite{metropolis1953} sought to calculate macroscopic thermodynamic quantities for systems of $N$ interacting particles in $d$ dimensions ($d \in \{2, 3\}$).

Let $x = (x_1, \dots, x_N) \in \mathcal{X}$ denote a spatial configuration with potential energy $E(x)$. At thermodynamic equilibrium, states follow the canonical Boltzmann distribution:
\begin{equation}
    \pi(x) = \frac{\widetilde{\pi}(x)}{Z} = \frac{\exp\left\{-\beta E(x)\right\}}{Z}, \qquad \text{where } \beta = \frac{1}{k_{\mathrm{B}}T}.
\end{equation}
Evaluating expected physical properties $\mathbb{E}_{\pi}[h(X)]$ required computing the ratio of two high-dimensional integrals over the partition function $Z = \int_{\mathcal{X}} \exp\{-\beta E(x)\} dx$. Standard Monte Carlo sampling $X_i \sim \text{Uniform}(\mathcal{X})$ failed because uniform samples in dense states almost always placed particles unrealistically close together, driving $E(x) \to \infty$ and yielding negligible weights $\exp\{-\beta E(x)\} \approx 0$.

To bypass evaluating $Z$, Metropolis et al. (1953) \cite{metropolis1953} introduced a guided random walk using a symmetric proposal distribution $q(y \mid x) = q(x \mid y)$. While published under five authors, historical accounts attribute the mathematical derivation to Marshall Rosenbluth and the implementation on the MANIAC I computer to Arianna Rosenbluth \cite{gubernatis2005}. W. K. Hastings (1970) \cite{hastings1970} generalised this to asymmetric proposals $q(y \mid x) \neq q(x \mid y)$, transforming a physical simulation technique into a general mathematical tool for statistical sampling.

\begin{proposition}[Cancellation of Normalising Constants \cite{robert2004mcmc}]
\label{prop:cancellation-constant}
Let $\widetilde{\pi}(\theta) = p(y \mid \theta)\pi_0(\theta)$ be an unnormalised target density with normalising constant $Z = p(y)$. For any proposal $q(\cdot \mid \theta)$, the Metropolis--Hastings acceptance probability $\alpha_{\mathrm{MH}}(\theta, \theta')$ satisfies:
\begin{equation}
    \alpha_{\mathrm{MH}}(\theta, \theta') = \min\left\{1, \; \frac{\pi(\theta') q(\theta \mid \theta')}{\pi(\theta) q(\theta' \mid \theta)}\right\} = \min\left\{1, \; \frac{\widetilde{\pi}(\theta') q(\theta \mid \theta')}{\widetilde{\pi}(\theta) q(\theta' \mid \theta)}\right\}.
\end{equation}
\end{proposition}

\noindent This cancellation is the foundation of MCMC computation. It allows samples to be drawn from complex target posterior distributions knowing only the unnormalised joint density $p(y \mid \theta)\pi_0(\theta)$, entirely bypassing the intractable high-dimensional integration required to evaluate $p(y)$.

\begin{theorem}[Detailed Balance and Ergodicity \cite{meyn2012mcke, tierney1994}]
\label{thm:mcmc-ergodic}
If a transition kernel $P(x, dy)$ satisfies detailed balance $\pi(dx) P(x, dy) = \pi(dy) P(y, dx)$, then $\pi$ is invariant. Furthermore, if the resulting Markov chain is $\pi$-irreducible, aperiodic, and positive Harris recurrent, then for any $\pi$-integrable function $h$:
\begin{equation}
    \lim_{M \to \infty} \frac{1}{M} \sum_{m=1}^{M} h(X_m) = \mathbb{E}_\pi[h(X)] \quad \text{a.s.}
\end{equation}
\end{theorem}

\noindent Detailed balance theorem offers a tractable, local condition to ensure that the chain preserves the target distribution $\pi$. Ergodicity provides the formal theoretical guarantee that sample path averages will almost surely converge to true posterior expectations as the chain runs infinitely long.

% -----------------------------------------------------------------------------
\subsubsection{Image Restoration and Gibbs Sampling (1984)}
\label{subsec:gibbs-sampling}

Independently of Hastings' statistical formulation, Stuart and Donald Geman (1984) \cite{geman1984} addressed the problem of recovering a clean digital image $F = (F_s)_{s \in \mathcal{S}}$ over a grid $\mathcal{S}$ given noisy observations $G$, via the posterior $\pi(f \mid g) \propto p(g \mid f) \pi_0(f)$.

To capture local spatial correlations, Geman and Geman modeled the prior $\pi_0(f)$ as a Markov Random Field (MRF) over local neighbourhoods $\partial s$.

\begin{theorem}[Hammersley--Clifford Theorem \cite{besag1974spatial, geman1984}]
\label{thm:hammersley-clifford}
A strictly positive random field $F$ on a finite lattice $\mathcal{S}$ is a Markov Random Field if and only if its joint distribution $\pi_0(f)$ is a Gibbs distribution:
\begin{equation}
    \pi_0(f) = \frac{1}{Z} \exp\left\{ -\frac{1}{T} \sum_{C \in \mathcal{C}} V_C(f) \right\},
\end{equation}
where $T > 0$ is temperature, $\mathcal{C}$ is the collection of local pixel cliques, and $V_C(f)$ are clique potential functions.
\end{theorem}

\noindent This theorem builds a crucial theoretical bridge between local spatial specifications and global joint probability fields. It proves that specifying simple, local interactions is mathematically equivalent to defining a valid joint distribution over the entire lattice.

For a lattice of $d$ pixels with $K$ discrete intensity levels, $Z$ involves $K^d$ terms (e.g., $256^{512 \times 512}$), rendering direct evaluation impossible. Geman and Geman introduced what they named the \emph{Gibbs sampler}—in honor of Josiah Willard Gibbs—noting that the full conditional distribution of a pixel $F_s$ given $F_{-s}$ depends solely on its local neighbourhood $\partial s$:
\begin{equation}
    \pi(f_s \mid f_{-s}) = \pi(f_s \mid f_{\partial s}) = \frac{\exp\left\{ -\frac{1}{T} \sum_{C : s \in C} V_C(f) \right\}}{\sum_{u \in \mathcal{F}_s} \exp\left\{ -\frac{1}{T} \sum_{C : s \in C} V_C(u, f_{-s}) \right\}},
\end{equation}


where $F_{-s}$ denotes the values of all pixels on the grid except pixel $s$. By updating each coordinate iteratively from its full conditional distribution, the chain targets the full joint distribution:
\begin{align}
    X_1^{(m+1)} &\sim \pi(x_1 \mid X_2^{(m)}, X_3^{(m)}, \dots, X_d^{(m)}), \\
    &\;\;\vdots \notag \\
    X_d^{(m+1)} &\sim \pi(x_d \mid X_1^{(m+1)}, X_2^{(m+1)}, \dots, X_{d-1}^{(m+1)}).
\end{align}


\begin{proposition}[Gibbs Update as Exact Metropolis--Hastings \cite{robert2004mcmc}]
\label{prop:gibbs-as-mh}
A Gibbs update proposing $y = (y_i, x_{-i}) \sim \pi(y_i \mid x_{-i})$ is equivalent to a Metropolis--Hastings step with acceptance probability identically equal to 1:
\begin{equation}
    \alpha(x,y) = \min\left\{1, \; \frac{\pi(y) q(x \mid y)}{\pi(x) q(y \mid x)}\right\} = \min\left\{1, \; \frac{\pi(x_{-i})\pi(y_i \mid x_{-i}) \cdot \pi(x_i \mid x_{-i})}{\pi(x_{-i})\pi(x_i \mid x_{-i}) \cdot \pi(y_i \mid x_{-i})}\right\} = 1.
\end{equation}
\end{proposition}
\noindent This unifies Gibbs sampling within the broader Metropolis--Hastings framework. It shows that sampling directly from full conditionals avoids rejected moves entirely ($\alpha = 1$), removing the need to manually choose and tune proposal densities.

% -----------------------------------------------------------------------------
\subsubsection{The Mainstream Statistical Revolution and Applications (1987--Present)}
\label{subsec:mcmc-applications}

The bridge between physical/image models and mainstream Bayesian statistics was completed by Tanner and Wong (1987) \cite{tannerwong1987} via \emph{Data Augmentation} for latent variables, and synthesized by Gelfand and Smith (1990) \cite{gelfandsmith1990}. Gelfand and Smith demonstrated that Gibbs sampling could routinely calculate marginal posterior distributions across arbitrary hierarchical Bayesian models.

The subsequent release of declarative MCMC engines like BUGS (Bayesian inference Using Gibbs Sampling) language \cite{gilks1994} laid the foundation for modern probabilistic programming by demonstrating how to specify complex hierarchical statistical models using Markov chain Monte Carlo (MCMC) methods. Key application domains include:

\begin{itemize}
    \item \textbf{Hierarchical Bayesian Models:} Sequential sampling of hyper-parameters and group parameters in multi-level structures \cite{gelman2013bayesian}.
    \item \textbf{Latent Variable \& Missing Data Models:} Treating unobserved states $Z$ as parameters via Data Augmentation to sample from $\pi(\theta \mid z, y)$ and $\pi(z \mid \theta, y)$ \cite{tannerwong1987}.
    \item \textbf{Spatial Statistics \& Image Processing:} Estimating spatial autoregressions and Markov Random Fields in disease mapping and geostatistics \cite{besag1974spatial}.
    \item \textbf{Phylogenetics \& Population Genetics:} Traversing tree topologies and estimating evolutionary rates from genomic sequences \cite{yang2014molecular}.
\end{itemize}



% End, updated 21-07-2026 -----------------------------------------------------------------------------

\subsection{Bayes through time: state-space models, Kalman filtering and particle filters}


Kalman Filters developed for Apollo Space mission:  applied in Apollo navigation.
https://www.youtube.com/watch?v=MMFpaNm-1EI




Standard MCMC assumes a \emph{static} posterior distribution $\pi(\theta \mid y)$. However, when observations arrive sequentially over time ($y_{1:t}$), re-running MCMC at each step $t$ is computationally prohibitive. 

\cite{sequentialMCinvitation2022}

This necessitates \emph{sequential Bayesian filtering}: evaluating $p(x_t \mid y_{1:t})$ recursively via prediction and update steps. The following section explores how exact updates (Kalman Filters) and particle approximations (Sequential Monte Carlo) extend Bayesian inference into time-evolving dynamical systems.

% \subsection{The path from distributions to optimisation: variational methods}

% - women, other countries

% \subsection{Bayesian Statistics powers Generative Language Models}

\newpage

% % --- START OF FIT-TO-PAGE DECISION TREE ---
% \begin{figure}[htbp]
% \centering

% % Color Palette
% \definecolor{primary}{HTML}{1A365D}    % Deep Blue
% \definecolor{decision}{HTML}{2C5282}   % Slate Blue
% \definecolor{leaf}{HTML}{276749}       % Forest Green

% % Guaranteed page-width fit
% \resizebox{\linewidth}{!}{%
% \begin{tikzpicture}[
%     auto,
%     node distance = 0.9cm and 0.8cm,
%     font = \sffamily\small,
%     % Decision Box (Rhombus/Diamond)
%     question/.style = {
%         diamond,
%         draw=decision,
%         fill=decision!12,
%         thick,
%         text width=2.8cm,
%         align=center,
%         inner sep=0pt,
%         drop shadow={opacity=0.15, shadow xshift=1pt, shadow yshift=-1pt}
%     },
%     % Algorithm/Leaf Box (Rectangle)
%     leaf/.style = {
%         rectangle,
%         rounded corners=3pt,
%         draw=leaf,
%         fill=leaf!12,
%         thick,
%         text width=3.2cm,
%         align=center,
%         minimum height=0.7cm,
%         inner sep=3pt,
%         drop shadow={opacity=0.15, shadow xshift=1pt, shadow yshift=-1pt}
%     },
%     branch/.style = {
%         font=\sffamily\bfseries\scriptsize,
%         color=primary
%     },
%     line/.style = {
%         draw=primary,
%         -Stealth,
%         thick
%     }
% ]

% % ==================== NODES ====================

% % Level 1: Q1 Likelihood
% \node [question] (q1) {\textbf{Q1: Likelihood}\\Is $p(x|z)$ computable?};
% \node [leaf, left=0.8cm of q1] (lfi) {\textbf{Likelihood-Free Inference}\\Classical ABC / Neural SBI};

% % Level 2: Q2 Data Type
% \node [question, below=0.8cm of q1] (q2) {\textbf{Q2: Data Type}\\Dynamic / streaming?};

% % Branch Q2 -> YES (Streaming Sub-tree)
% \node [question, right=0.8cm of q2] (q2a) {\textbf{Q2a: State}\\Continuous or Discrete?};
% \node [leaf, above right=0.1cm and 0.8cm of q2a] (hmm) {\textbf{HMM}\\Forward-Backward / Viterbi};

% \node [question, below=0.7cm of q2a] (q2b) {\textbf{Q2b: Dynamics}\\Linear \& Gaussian?};
% \node [leaf, right=0.8cm of q2b] (kf) {\textbf{Kalman Filter (KF)}};

% \node [question, below=0.7cm of q2b] (q2c) {\textbf{Q2c: Dynamics}\\Severity of non-linearity?};
% \node [leaf, right=0.8cm of q2c] (ekf) {\textbf{EKF / UKF}};

% \node [question, below=0.7cm of q2c] (q2d) {\textbf{Q2d: Parameters}\\Need static params ($\theta$)?};
% \node [leaf, right=0.8cm of q2d] (smc) {\textbf{SIR Particle Filter}};
% \node [leaf, below=0.3cm of smc] (pmcmc) {\textbf{PMCMC / SMC$^2$}};

% % Branch Q2 -> NO (Static Data)
% \node [question, below=3.8cm of q2] (q3) {\textbf{Q3: Data Scale}\\Massive dataset ($N > 10^5$)?};

% % Branch Q3 -> YES (Massive N)
% \node [question, right=0.8cm of q3] (q3a) {\textbf{Q3a: Latents}\\High-D continuous generative model?};
% \node [leaf, above right=0.1cm and 0.8cm of q3a] (vae) {\textbf{Amortized VI (VAE)}};
% \node [leaf, below right=0.1cm and 0.8cm of q3a] (svi) {\textbf{Stochastic VI / SGLD}};

% % Branch Q3 -> NO (Moderate/Small Data)
% \node [question, below=1.8cm of q3] (q4) {\textbf{Q4: Parameters}\\Continuous or Discrete?};

% % Branch Q4 -> Discrete
% \node [question, right=0.8cm of q4] (q4a) {\textbf{Q4a: Target}\\Multimodal / Combinatorial?};
% \node [leaf, above right=0.1cm and 0.8cm of q4a] (asmc) {\textbf{Annealed SMC}};
% \node [leaf, below right=0.1cm and 0.8cm of q4a] (gibbs) {\textbf{Gibbs / MH}};

% % Branch Q4 -> Continuous
% \node [question, below=1.6cm of q4] (q4b) {\textbf{Q4b: Dimension}\\High-D ($D > 20$) or Correlated?};
% \node [leaf, right=0.8cm of q4b] (hmc) {\textbf{HMC / NUTS}\\($\nabla_z \log p(x,z)$)};
% \node [leaf, below=0.3cm of hmc] (rwm) {\textbf{Random-Walk MH / Slice}};


% % ==================== PATHS & CONNECTIONS ====================

% % Main Spine
% \draw [line] (q1) -- node [branch, left] {NO} (lfi);
% \draw [line] (q1) -- node [branch, right] {YES} (q2);
% \draw [line] (q2) -- node [branch, right] {NO (Static)} (q3);
% \draw [line] (q3) -- node [branch, right] {NO} (q4);

% % Dynamic Branch (Q2)
% \draw [line] (q2) -- node [branch, above] {YES} (q2a);
% \draw [line] (q2a) |- node [branch, near start, above] {Discrete} (hmm);
% \draw [line] (q2a) -- node [branch, right] {Continuous} (q2b);
% \draw [line] (q2b) -- node [branch, above] {YES} (kf);
% \draw [line] (q2b) -- node [branch, right] {NO} (q2c);
% \draw [line] (q2c) -- node [branch, above] {Mild} (ekf);
% \draw [line] (q2c) -- node [branch, right] {Severe} (q2d);
% \draw [line] (q2d) -- node [branch, above] {NO} (smc);
% \draw [line] (q2d) |- node [branch, near start, right] {YES} (pmcmc);

% % Massive Scale Branch (Q3)
% \draw [line] (q3) -- node [branch, above] {YES} (q3a);
% \draw [line] (q3a) |- node [branch, near start, above] {YES} (vae);
% \draw [line] (q3a) |- node [branch, near start, below] {NO} (svi);

% % Parameter Space Branches (Q4)
% \draw [line] (q4) -- node [branch, above] {Discrete / Mixed} (q4a);
% \draw [line] (q4) -- node [branch, right] {Continuous} (q4b);

% \draw [line] (q4a) |- node [branch, near start, above] {YES} (asmc);
% \draw [line] (q4a) |- node [branch, near start, below] {NO} (gibbs);

% \draw [line] (q4b) -- node [branch, above] {YES} (hmc);
% \draw [line] (q4b) |- node [branch, near start, right] {NO} (rwm);

% \end{tikzpicture}%
% }
% \caption{Master Decision Tree for Selecting Bayesian Inference Methods.}
% \label{fig:bayes_decision_tree}
% \end{figure}
% --- END OF FIT-TO-PAGE DECISION TREE ---

\newpage



% \subsection{Dark side of Bayes Theorem vs additional Bayes developments?}

% Sally Clark 
% https://theconversation.com/bayes-theorem-the-maths-tool-we-probably-use-every-day-but-what-is-it-76140


% Data Privacy Issues - inverse engineer NHS data.
% -- cases where somehting has happened 

% neuclear weapon from Neumann.. 

% medical paradox - https://www.youtube.com/watch?v=lG4VkPoG3ko 
% % Can remove if not including this section
% \begin{figure}
%     \centering
%     \includegraphics[width=0.5\linewidth]{Medical paradox.png}
%     \caption{Medical Paradox}
%     \label{fig:placeholder}
% \end{figure}





